% Copyright (c) 2026 Eugene Joseph M. Ragasa.
% SPDX-License-Identifier: Apache-2.0
% Author: Eugene Joseph M. Ragasa
% Formal proof completed: 2026-10-19
\documentclass[11pt]{article}
\usepackage[T1]{fontenc}
\usepackage{amsmath,amssymb,booktabs}
\usepackage[margin=1in]{geometry}
\usepackage[hidelinks]{hyperref}
\usepackage{tabularx}
\title{Proof Dossier: Traceless--Identity Frobenius Orthogonality}
\author{Eugene Joseph M. Ragasa}
\date{Formal proof completed: 2026-10-19}
\begin{document}
\maketitle
\section{Identity and status}
\begin{tabularx}{\textwidth}{@{}l>{\raggedright\arraybackslash}X@{}}
\toprule
Statement ID & \texttt{BG-HILBERT-THM-010} \\
Status & \texttt{proved} \\
Lean declaration & \path{ProjectKoios.BandGauge.Hilbert.rankTwoWeightedFrobeniusCorrelation_identity_eq_zero} \\
Module & \path{ProjectKoios/BandGauge/Hilbert/WeightedFrobenius.lean} \\
\bottomrule
\end{tabularx}

\section{Exact statement}
If every matrix $X_i$ in a finite rank-two field has trace zero, then its weighted real
Frobenius correlation with the constant identity field vanishes:
\[
 C(I,X)=0.
\]
No weight sign or normalization assumption is needed.

\section{Source and proof}
This theorem \texttt{formalizes} the vanishing shift cross term in Supplementary Section
S1.1 of Ragasa2026A. Owner: \path{ksdft2effmass}; immutable commit
\begin{center}\small\texttt{99a46cb59d0c6f27ada7d709f194996ef9619d9f}.\end{center}
Source: \path{docs/publications/papers/Ragasa2026a/v3/}\path{supplementary.tex},
lines 69--76. At each coordinate, the Frobenius pairing with identity is the real part of
the trace. Pointwise trace zero makes every weighted summand zero.

\section{Boundary and trusted base}
The theorem assumes tracelessness; it does not prove it for external matrices. Lean and
pinned Mathlib report \texttt{propext}, \texttt{Classical.choice}, and
\texttt{Quot.sound}. No admission or scientific axiom is used. The proof and prose were
AI-assisted and kernel-checked. No independent expert review is recorded. Status:
\texttt{proved}; working draft.
\end{document}
